% 
%  chapter8.tex
%  ISELthesis
%  
%  Created by Matilde Pós-de-Mina Pato on 2012/10/09.
%
\chapter{Conclusions and Future Work}
\label{cha:conclusions}

This dissertation set out to unify two previously independent tools --- QuickFaaS, which deploys
serverless functions, and OmniFlow, which defines and deploys workflows --- so that a serverless
application can be deployed from a single workflow definition, without manually coordinating
function endpoints. This final chapter summarises the contribution, revisits the objectives,
assesses the work critically, and outlines future directions.

\section{Summary of Contributions}
\label{sec:concl-contributions}

The work delivered two complementary halves.

\paragraph{Unification (core).}
OmniFlow was extended with a \emph{Function Registry} and a deployment-time \emph{resolution
cascade}. Internal workflow calls declare a stable logical identifier through the
\texttt{internalFunction(name, deploymentDescriptorPath)} construct added to the \texttt{call}
step, mutually exclusive with the literal \texttt{host}/\texttt{path} of external calls. At
deployment time OmniFlow traverses the workflow tree, and for each internal call resolves the
endpoint from the registry, discovers it directly on the provider, or --- as a last resort ---
deploys it through QuickFaaS, writing the resulting metadata back to the registry. Function
deployment is abstracted behind the \texttt{InternalFunctionDeployer} strategy interface, which
keeps the two providers symmetric and makes support for a further platform a matter of adding one
resolver and one deployer.

\paragraph{AWS support (complement).}
QuickFaaS, previously limited to GCP and Azure, was extended with an AWS provider, and OmniFlow
gained an \texttt{AwsLambdaDeployer} that performs the AWS-specific work (IAM execution role,
region resolution from the S3 bucket, readiness polling, and granting Step Functions permission to
invoke the Lambda). Internal AWS calls are rendered with the native
\texttt{arn:aws:states:::lambda:invoke} integration rather than an HTTP task, keeping the
invocation within the AWS control plane.

\section{Objectives Revisited}
\label{sec:concl-objectives}

The two research questions of Section~\ref{sec:int_objectives} are answered as follows:
\begin{itemize}
  \item \textbf{RQ1 (feasibility and correctness)} --- yes, for the resolution logic. On both AWS
  and GCP the cascade binds each internal call to the right endpoint and deploys a function only
  when its absence is confirmed, and it aborts on stale entries, ambiguous names and inaccessible
  regions (Section~\ref{sec:eval-correctness}). The evidence comes from unit tests that simulate
  the provider, and on GCP bindings to first-generation Cloud Functions are not validated live.
  \item \textbf{RQ2 (overhead)} --- small. In the production resolvers, resolving a workflow with
  up to 200 internal calls takes at most a few milliseconds locally, growing linearly with the
  number of internal calls, while the registry is read once per deployment
  (Section~\ref{sec:eval-discussion}). These figures exclude the provider calls that validation,
  discovery and deployment make.
\end{itemize}

Against the four objectives stated in Section~\ref{sec:int_objectives}, the work stands as
follows:
\begin{enumerate}
  \item \textbf{Design a unified integration architecture} --- \emph{met.} The registry, the
  \texttt{functionRef} indirection, the \texttt{internalFunction} construct and the resolver
  strategy (Chapters~\ref{ch:proposed_solution} and~\ref{cha:impl}) let a single workflow
  definition drive both function and workflow deployment, binding endpoints when the workflow is
  deployed (\S\ref{sec:resolution-cascade}).
  \item \textbf{Implement an orchestration prototype} --- \emph{met on AWS, partially met on GCP.}
  The prototype resolves the internal functions of a workflow on both providers, and deploys the
  missing ones, without endpoints being copied by hand. On GCP, however, the first-generation Cloud
  Functions that QuickFaaS deploys are not covered by validation or discovery: a stale binding to
  one of them goes undetected, and a binding lost from the registry cannot be rediscovered
  (\S\ref{subsec:store-internals}).
  \item \textbf{Extend QuickFaaS to AWS Lambda} --- \emph{met.} QuickFaaS gained an AWS provider
  and OmniFlow an \texttt{AwsLambdaDeployer} (Sections~\ref{subsec:quickfaas-aws-provider}
  and~\ref{subsec:omniflow-aws-deployer}). The provider-agnostic packaging layer adds about
  0.9\,KB to a Lambda bundle (Table~\ref{tab:eval-s1}), and the resolver tests cover both
  providers (Section~\ref{sec:eval-correctness}).
  \item \textbf{Evaluate feasibility and overhead} --- \emph{partially met.}
  Chapter~\ref{cha:evaluation} checks the decisions of the resolution cascade on both providers
  and measures the resolution overhead, the registry write path and the packaging size. The
  evidence is local only: the correctness checks simulate the provider, the benchmarks ran with a
  short configuration on a shared machine, and neither end-to-end deployment times nor the
  provider calls that validation and discovery make were measured.
\end{enumerate}

\section{Critical Assessment}
\label{sec:concl-critical}

The implementation is that of a prototype, and several limitations should be stated plainly.

\begin{itemize}
  \item \textbf{Static binding.} Endpoints are resolved when the workflow is deployed and embedded
  in the rendered artifact (\S\ref{sec:resolution-cascade}). An endpoint that changes afterwards
  reaches a deployed workflow only when that workflow is deployed again; nothing resolves
  endpoints at run time.
  \item \textbf{No update path.} The cascade deploys a function only when it is missing and never
  updates one that exists: a change to a function's code is not deployed by a workflow deployment,
  and has to be deployed separately (for example, with QuickFaaS directly).
  \item \textbf{Unguarded registry.} The registry is a flat JSON file with no cache, no integrity
  or authenticity guarantee, and no cross-process locking; concurrent writers race, and a tampered
  \texttt{url}/ARN would be trusted blindly.
  \item \textbf{Partial GCP coverage.} On GCP the resolver validates and discovers Cloud Run
  services only, while QuickFaaS deploys first-generation Cloud Functions. Bindings to those
  functions are trusted without a live check, and a binding lost from the registry cannot be
  rediscovered (\S\ref{subsec:store-internals}). The AWS path provides both guarantees.
  \item \textbf{Broad, automatic IAM.} The AWS path auto-creates and mutates a shared execution
  role and grants invoke permissions as a deployment side effect --- convenient for a demo, but at
  odds with a least-privilege, change-controlled environment.
  \item \textbf{Brittle integration boundary.} QuickFaaS is invoked as a subprocess, coupling two
  build systems across a process boundary and keying error handling off child-process text; the
  rationale for this boundary, and the in-process alternative, is discussed in
  Section~\ref{sec:design-rationale}.
  \item \textbf{Testability of cloud paths.} By design, the classes that call the cloud SDKs have
  no automated coverage, so the integration edges --- where production incidents occur --- are the
  least tested.
\end{itemize}

\section{Applicability}
\label{sec:concl-applicability}

The shape of the solution --- logical indirection between workflows and endpoints, provider
symmetry, and explicit fail-fast diagnostics --- fits the needs of an organisation that deploys the same
orchestration to several accounts or projects, and that must be able to move it to another provider
(for example, to meet the exit-strategy requirements of the EU Digital Operational Resilience
Act~\cite{dora2022}, or to expose open-banking functions under PSD2). Reaching that
setting would require hardening the gaps above: a signed, access-controlled registry backend, an
attributable write history for audit, and least-privilege IAM under change control. None of these
changes the resolution cascade itself; they harden the parts around it.
Chapter~\ref{cha:case-study} works through a concrete instance of this setting --- a card-payment
fraud-screening workflow --- to illustrate both the value of the integration and the extent of
this hardening.

\section{Future Work}
\label{sec:concl-future}

The following directions are ordered from smallest change with the clearest payoff to the larger
structural ones.

\begin{enumerate}
  \item \textbf{Cover first-generation Cloud Functions on GCP} with a Cloud Functions~v1 client
  used for validation, discovery and bootstrap, giving the GCP path the same drift detection and
  rediscovery guarantees as the AWS path.
  \item \textbf{Add an update path:} record a hash of each function's source in its registry entry
  and redeploy the function through QuickFaaS when the source its descriptor points to changes, so
  that a workflow deployment also brings its functions up to date.
  \item \textbf{Harden the registry:} a pluggable backend (database or secret/parameter store)
  behind the current interface, with cross-process locking, per-entry integrity, an in-memory
  cache with explicit invalidation, and an append-only, attributable write history.
  \item \textbf{Replace the subprocess boundary} by invoking QuickFaaS in-process as a library,
  removing the two-build-system coupling and the fragile parsing of subprocess output.
  \item \textbf{Move to least-privilege, change-controlled IAM,} with pre-provisioned roles and a
  dry-run mode that reports the IAM changes it would make instead of applying them silently.
  \item \textbf{Introduce test seams} for the cloud-calling code, extracting the SDK interactions
  behind narrow interfaces so the deploy/IAM/readiness logic can be unit-tested against fakes.
  \item \textbf{Extend the provider set,} adding Azure Functions (already a QuickFaaS target)
  through the existing extension point to further substantiate the portability claim.
\end{enumerate}

\section{Final Remarks}
\label{sec:concl-final}

The integration shows that function deployment and workflow orchestration can be driven from a
single, provider-agnostic definition, with concrete endpoint binding deferred to deployment time
through an explicit resolution cascade. In local measurements the overhead of this unification is
small and, with the single-read optimization applied in both production resolvers, grows linearly.
What remains is evidence against the real providers and engineering maturity --- registry
integrity, least-privilege operation, test coverage of the cloud edges and an update path --- to
take the prototype to production.
